\chapter{Methodology}
\label{ch:methodology}

\section{Chapter Overview}
\label{sec:meth_chapter_overview}
Building a customizable paper virtual keyboard using only a normal webcam and deep learning requires a clear and systematic engineering methodology. Unlike regular mechanical keyboards or touchscreens, a printed sheet of paper has no physical buttons, key switches, or electrical wiring. Instead, user typing actions are detected using computer vision, perspective homography, scale-normalized hand tracking, and neural network sequence classification running in real time on standard computer processors without requiring a GPU.

This chapter explains the research framework, system design, and software pipelines developed for this project. Section~\ref{sec:meth_framework} introduces the scientific research paradigm and the six stages of the Design Science Research Methodology (DSRM). Section~\ref{sec:meth_architecture} outlines the overall software architecture and design patterns connecting the two main applications. Section~\ref{sec:meth_layout_designer} describes Application 1 (Layout Designer), covering key placement, border AprilTag anchors, and file export. Section~\ref{sec:meth_data_pipeline} details the dataset collection, 12~FPS video resampling, 21-joint MediaPipe hand tracking, scale normalization, and sliding window creation. Section~\ref{sec:meth_dl_models} presents the multi-finger touch detection problem, loss functions, and the PyTorch LSTM model. Section~\ref{sec:meth_runtime_engine} explains Application 2 (Runtime Virtual Keyboard Engine), covering planar homography, coordinate transformation, and action multiplexing. Section~\ref{sec:meth_summary} summarizes the chapter.

\section{Research Framework and Scientific Methodology}
\label{sec:meth_framework}

\subsection{Research Paradigm and Deductive Reasoning}
\label{subsec:meth_paradigm}
This research follows an \textbf{Empirical Positivist and Pragmatist Paradigm} with a \textbf{Deductive Logical Approach}. In computer vision and Human-Computer Interaction (HCI), finger movements, camera perspectives, homography matrices, and impact decelerations are observable physical and mathematical events. A positivist approach allows us to measure these phenomena using clear quantitative numbers, including precision, recall, F1-score, reprojection error in millimeters, and runtime speed in milliseconds, across repeatable laboratory tests. Pragmatism further supports our engineering decisions by prioritizing practical, real-world utility: achieving real-time interaction on commodity CPUs without requiring expensive GPUs or specialized depth cameras.

We used deductive reasoning to apply established scientific concepts to our software design:
\begin{enumerate}
    \item \textbf{Planar Homography Theory:} From projective geometry \cite{Pirchheim2011Homography}, four coplanar reference points on a sheet of paper are mathematically sufficient to calculate a unique $3 \times 3$ homography matrix $H$. This matrix corrects for camera tilt angles and paper movement on a desk.
    \item \textbf{Scale Invariance Under Perspective Projection:} In camera optics, the perceived size of an object changes inversely with its distance from the camera lens. By dividing hand landmark coordinates by the measured anatomical hand length ($L_{\text{hand}}$), distance and camera resolution differences are cancelled out before classification.
    \item \textbf{Kinematic Deceleration During Physical Contact:} When a typing finger hits a solid desk through a sheet of paper, its downward speed drops sharply to zero. Detecting this rapid deceleration across temporal windows allows the system to recognize physical touch contact without needing specialized 3D depth sensors.
\end{enumerate}

\subsection{Methodological Grounding via Saunders' Research Onion}
\label{subsec:meth_saunders_onion}
To systematically structure our research process from philosophical assumptions to practical data analysis, the methodology aligns with the six layers of Saunders' Research Onion \cite{ReviewVirtualKeyboard2020}:
\begin{itemize}
    \item \textbf{Layer 1 (Research Philosophy):} Positivism and Pragmatism, focusing on objective empirical measurements (latency, F1-scores, character error rates) and real-world system utility on low-cost hardware.
    \item \textbf{Layer 2 (Research Approach):} Deductive reasoning, developing mathematical models of planar perspective transformation and temporal deceleration, and experimentally testing them against empirical data.
    \item \textbf{Layer 3 (Research Strategy):} Design Science Research Methodology (DSRM) combined with experimental laboratory benchmarking across modular software artifacts.
    \item \textbf{Layer 4 (Methodological Choice):} Dominantly quantitative empirical research, evaluating model classification metrics and system latency, supplemented by standardized user experience evaluation.
    \item \textbf{Layer 5 (Time Horizon):} Cross-sectional study design, conducting comprehensive performance and user evaluations across distinct camera angles, lighting conditions, and resolution levels.
    \item \textbf{Layer 6 (Techniques and Procedures):} 12~FPS video capture, 21-landmark skeletal tracking, unitless hand-length normalization, 5-frame temporal windowing, PyTorch LSTM training, and planar homography transformation.
\end{itemize}

\subsection{Design Science Research Methodology (DSRM)}
\label{subsec:meth_dsrm}
The software engineering and development process followed the six stages of the Design Science Research Methodology (DSRM) \cite{ReviewVirtualKeyboard2020}, as shown in Figure~\ref{fig:dsrm_process_flow}.

\begin{figure}[htbp]
\centering
\resizebox{0.92\textwidth}{!}{%
\begin{tikzpicture}[
    phase/.style = {rectangle, draw=blue!80!black, fill=blue!10!white, thick, rounded corners=4pt, align=center, font=\sffamily\bfseries\footnotesize, text width=2.2cm, minimum height=1.3cm, inner sep=4pt},
    arrow/.style = {draw=blue!80!black, ->, >=Stealth, thick}
]
\node (p1) [phase] at (0,0) {Phase 1\\Problem\\Identification};
\node (p2) [phase] at (2.8,0) {Phase 2\\Define Objective\\of Solution};
\node (p3) [phase] at (5.6,0) {Phase 3\\Design and\\Development};
\node (p4) [phase] at (8.4,0) {Phase 4\\Demonstration};
\node (p5) [phase] at (11.2,0) {Phase 5\\Evaluation};
\node (p6) [phase] at (14.0,0) {Phase 6\\Communication};

\draw [arrow] (p1) -- (p2);
\draw [arrow] (p2) -- (p3);
\draw [arrow] (p3) -- (p4);
\draw [arrow] (p4) -- (p5);
\draw [arrow] (p5) -- (p6);
\draw [arrow] (p5.south) -- ++(0,-0.6) -| node[below,midway,font=\sffamily\scriptsize]{Iterative Algorithmic Refinement} (p3.south);
\end{tikzpicture}%
}
\caption{Design Science Research Methodology (DSRM) execution framework.}
\label{fig:dsrm_process_flow}
\end{figure}

The six DSRM stages were executed as follows:
\begin{itemize}
    \item \textbf{Phase 1: Problem Identification:} Identified key weaknesses in existing virtual keyboards from the literature, including dependence on expensive depth cameras, failures under room shadows, fixed camera setups, and single-finger tracking restrictions.
    \item \textbf{Phase 2: Objectives Formulation:} Formulated five quantitative technical goals, targeting real-time execution on common CPUs ($< 83.33\text{ ms}$ per frame at 12~FPS), touch detection F1-score above 0.94 across all five fingers, and decoupling paper geometry from software actions.
    \item \textbf{Phase 3: Design and Development:} Developed the software suite in two parts: Application 1 (Layout Designer) for creating printable layouts, and Application 2 (Runtime Engine) combining MediaPipe hand pose tracking, scale normalization, AprilTag homography, and PyTorch LSTM touch detection.
    \item \textbf{Phase 4: Demonstration:} Tested the system on ordinary desks using normal paper printed with an office laser printer and captured by a basic USB webcam.
    \item \textbf{Phase 5: Evaluation:} Conducted benchmark tests across 22 neural network models, measured latency across video resolutions from 360p to 1080p on CPU, and recorded typing speed (WPM) and error rates with real users.
    \item \textbf{Phase 6: Communication:} Documented the project rationale, system design, and experimental outcomes in this dissertation.
\end{itemize}

\subsection{Partitioned Development and Experimental Strategy}
\label{subsec:meth_partitions}
To make development and debugging manageable, the project was organized into five isolated technical partitions:
\begin{enumerate}
    \item \textbf{Partition 1 (Layout Design and Homography Simulation):} Built the PySide6 designer (\texttt{designer\_app.py}) and simulated touch coordinates in \texttt{designer/analyzer/} to verify that homography matrix calculations correctly map camera pixel positions to layout keys before building live camera code.
    \item \textbf{Partition 2 (Dataset Creation and Pipeline Engineering):} Implemented scripts in \texttt{mediapipeDetector/datacreator/} for video resampling at 12~FPS (\texttt{resample\_12fps.py}), 21-joint hand landmark extraction, unitless scale normalization (\texttt{normalize\_landmarks.py}), velocity calculations, and 5-frame sliding window generation.
    \item \textbf{Partition 3 (Deep Learning Model Benchmarking):} Evaluated 22 pure 2D model architecture combinations using \texttt{mediapipeDetector/deepLearningModels/run\_all.py} across five core model families (1D CNN, Attention, ResNet, BiLSTM, LSTM), selecting the optimal PyTorch touch detection model (\texttt{best\_finger\_touch\_lstm.pth}).
    \item \textbf{Partition 4 (Real-Time Responsiveness Evaluation):} Developed live camera testing scripts in \texttt{mediapipeDetector/realtimeprocess/} to verify CPU processing speed, multi-threaded windowing, and touch trigger stability.
    \item \textbf{Partition 5 (AprilTag Fiducial Tracking):} Implemented camera calibration and fiducial tracking scripts in \texttt{aprilTag/} to detect corner coordinates and calculate the $3 \times 3$ planar homography matrix $H$ under camera tilt angles up to $75^\circ$.
\end{enumerate}

\section{System Architecture and Software Design}
\label{sec:meth_architecture}
The virtual keyboard separates layout creation from live runtime processing. Figure~\ref{fig:system_architecture_framework} shows the system architecture connecting both applications.

\begin{figure}[htbp]
\centering
\resizebox{0.92\textwidth}{!}{%
\begin{tikzpicture}[
    node distance = 1.2cm and 1.5cm,
    appbox/.style = {rectangle, draw=blue!80!black, fill=blue!12!white, thick, rounded corners=5pt, align=center, font=\sffamily\bfseries\footnotesize, text width=13.0cm, inner sep=6pt},
    artbox/.style = {rectangle, draw=purple!80!black, fill=purple!10!white, thick, rounded corners=4pt, align=center, font=\sffamily\bfseries\scriptsize, text width=5.6cm, minimum height=1.0cm, inner sep=4pt},
    compbox/.style = {rectangle, draw=teal!80!black, fill=teal!10!white, thick, rounded corners=4pt, align=left, font=\sffamily\scriptsize, text width=5.8cm, minimum height=2.4cm, inner sep=6pt},
    arrow/.style = {draw=blue!80!black, ->, >=Stealth, thick}
]

\node (app1) [appbox] at (0, 0) {
\textbf{APPLICATION 1: LAYOUT DESIGNER SUITE (\texttt{designer/designer\_app.py})}\\
{\normalfont\sffamily\tiny Interactive PySide6 GUI Canvas: Drag-and-drop key placement, AprilTag anchor borders, and synchronized exports}
};

\node (pdf_art) [artbox] at (-3.6, -2.0) {Printable Layout PDF\\(Embedded AprilTag Anchors)};
\node (xml_art) [artbox] at (3.6, -2.0) {Layout XML Specification File\\(Bounding Boxes and Action Schemas)};

\node (app2) [appbox] at (0, -4.0) {
\textbf{APPLICATION 2: RUNTIME VIRTUAL KEYBOARD ENGINE}\\
{\normalfont\sffamily\tiny Setup GUI Configuration and Multi-Threaded Background Vision Execution Engine}
};

\node (comp_a) [compbox] at (-3.6, -6.6) {
\textbf{Component A: Setup and Mapper GUI}\\
- Select active monocular camera feed\\
- Load layout XML coordinate specification\\
- Interactive key-to-command mapping table\\
- Save / Load action configuration profiles (JSON)
};

\node (comp_b) [compbox] at (3.6, -6.6) {
\textbf{Component B: Background Runtime Engine}\\
- AprilTag tracker $\to$ Homography matrix $H$\\
- MediaPipe hand pose $\to$ Scale normalization ($L_{\text{hand}}$)\\
- 5-frame temporal window $\to$ PyTorch LSTM (CPU)\\
- Planar mapping: $P_{\text{XML}} = H \cdot P_{\text{pixel}}$\\
- Key lookup $\to$ Dispatch OS keystroke / script
};

\draw [arrow] (app1.south -| pdf_art.north) -- (pdf_art.north);
\draw [arrow] (app1.south -| xml_art.north) -- (xml_art.north);
\draw [arrow] (pdf_art.south) -- node[left, font=\sffamily\tiny, align=right]{Physical Print\ } (app2.north -| pdf_art.south);
\draw [arrow] (xml_art.south) -- node[right, font=\sffamily\tiny, align=left]{\ Loaded into App 2} (app2.north -| xml_art.south);
\draw [arrow] (app2.south -| comp_a.north) -- (comp_a.north);
\draw [arrow] (app2.south -| comp_b.north) -- (comp_b.north);

\end{tikzpicture}%
}
\caption{System architecture framework connecting Application 1 and Application 2 workflows.}
\label{fig:system_architecture_framework}
\end{figure}

To ensure stability, modularity, and responsiveness on standard computers, four standard software design patterns were applied:
\begin{enumerate}
    \item \textbf{Model-View-Controller (MVC):} In Application 1, the visual drawing canvas (View) is separated from key coordinate models (Model) and mouse drag events (Controller).
    \item \textbf{Producer-Consumer Multi-Threading:} In Application 2, webcam video acquisition runs in a separate thread (\texttt{camera\_thread.py}), putting frames into a queue. The vision processing engine (\texttt{model\_manager.py}) reads frames asynchronously, keeping the graphical interface responsive and preventing dropped frames.
    \item \textbf{Strategy Pattern for Action Dispatching:} Key press triggers are handled through distinct action classes, allowing easy switching between virtual OS keystrokes (using \texttt{pyautogui} or \texttt{pynput}) and system shell scripts (using Python \texttt{subprocess}).
    \item \textbf{Singleton Model Registry:} The PyTorch touch model weights (\texttt{best\_finger\_touch\_lstm.pth}) are loaded into CPU memory once during startup to avoid reloading delays while typing.
\end{enumerate}

\section{Application 1: Layout Designer Suite}
\label{sec:meth_layout_designer}
Application 1 (\texttt{designer/designer\_app.py}) provides an interactive PySide6 graphical program for designing custom keyboard layouts on an A4 page ($210 \times 297\text{ mm}$). Users can create, resize, snap, and position rectangular key buttons across a calibrated canvas grid.

\subsection{Key Geometry and Fiducial Anchor Positioning}
\label{subsec:meth_key_geometry}
Each key button $k$ has a unique key ID, a text label, and bounding box coordinates:
\begin{equation}
B_k = \left( x_{\min}^{(k)}, y_{\min}^{(k)}, x_{\max}^{(k)}, y_{\max}^{(k)} \right)
\label{eq:meth_key_box}
\end{equation}
where values are measured in millimeters relative to standard A4 sheet boundaries.

To enable camera tracking, the designer automatically embeds four AprilTag markers (from the \texttt{tag36h11} family, Tag IDs 0, 1, 2, and 3) along the page borders. Each marker has a fixed physical size ($s = 35\text{ mm}$) and known metric positions:
\begin{equation}
\mathbf{C}_{\text{tag}} = \left\{ (x_{\text{TL}}, y_{\text{TL}}), (x_{\text{TR}}, y_{\text{TR}}), (x_{\text{BR}}, y_{\text{BR}}), (x_{\text{BL}}, y_{\text{BL}}) \right\}
\label{eq:meth_tag_coords}
\end{equation}
Placing markers strictly along the outer borders keeps the central typing area open and prevents typing hands from covering all four tags at the same time.

\subsection{Synchronized Export Engine}
\label{subsec:meth_export_engine}
When a layout design is finished, the exporter creates two synchronized files:
\begin{enumerate}
    \item \textbf{Printable Vector PDF:} Rendered at 300 DPI containing key outlines, key labels, and high-contrast AprilTag markers. This sheet is printed on standard desktop paper.
    \item \textbf{Layout Specification XML:} An XML document defining paper size, AprilTag anchor positions, and bounding box dimensions for all keys. This file is loaded into Application 2.
\end{enumerate}

\section{Data Collection and Kinematic Preprocessing Pipeline}
\label{sec:meth_data_pipeline}
Detecting touch events accurately using an ordinary 2D webcam requires tracking finger movement over time. The feature pipeline converts incoming RGB video into normalized temporal feature vectors.

\subsection{Dataset Collection Protocol and Methodological Rigor}
\label{subsec:meth_dataset_capture}
To train and benchmark the neural network, typing video recordings were captured under controlled laboratory conditions following strict empirical rigor:
\begin{itemize}
    \item \textbf{Camera Setup:} Standard 1080p monocular webcams placed between $35\text{ cm}$ and $65\text{ cm}$ above the desk, with viewing angles from $30^\circ$ to $75^\circ$ relative to the surface plane.
    \item \textbf{Lighting Environments:} Three different room lighting levels were tested: Dim Room ($150\text{ lux}$), Standard Office ($500\text{ lux}$), and Bright Light ($2,500+\text{ lux}$).
    \item \textbf{Participants:} 15 university participants (ages 20 to 32) with varying hand sizes ($L_{\text{hand}}$ from $14.2\text{ cm}$ to $21.5\text{ cm}$), diverse skin tones across the Fitzpatrick scale (Types I to VI), and different typing speeds.
    \item \textbf{Typing Actions:} Participants performed single-finger key taps, multi-finger chords, hand resting poses, hovering motions, and regular typing across all five digits.
\end{itemize}

Ground truth touch timestamps were collected using small low-voltage electrical contact sensors attached to the paper surface, providing accurate contact labels $y_t \in \{0, 1\}^5$ aligned with each video frame.

\textbf{Validity and Reliability Assurance:} In accordance with empirical computing research principles:
\begin{itemize}
    \item \textbf{Internal Validity:} Confounds were controlled by maintaining identical camera sensor exposure settings, fixed 12~FPS timing intervals, and synchronized hardware timestamps during recording sessions.
    \item \textbf{Construct Validity:} Ground truth contact was verified using physical electrical contact signals rather than subjective human video annotation, ensuring that contact labels reflect physical paper impact.
    \item \textbf{External Validity (Generalisability):} Diversity in participant hand sizes, skin tones, ambient illuminations, and camera angles ensured that trained models generalize across real-world desktop environments.
    \item \textbf{Triangulation:} The research incorporates data triangulation (simulated touch datasets, empirical video datasets, and live user typing trials) and methodological triangulation (offline cross-validation across 5 deep learning families and live CPU runtime profiling).
\end{itemize}

\textbf{Research Ethics and Participant Consent:} In line with ethical research standards, all 15 participants provided informed written consent before participating in typing sessions. Data collection was strictly limited to fingertip and hand coordinate positions; no facial images, personal identifiers, or sensitive demographic information were stored. Participants were informed of their right to withdraw at any stage without consequence.

\subsection{12 FPS Sub-Sampling and Single-Hand Tracking}
\label{subsec:meth_subsampling}
Camera video streams are sampled at a uniform rate of \textbf{12 Frames Per Second (FPS)} ($\Delta t = \frac{1}{12}\text{ s} \approx 83.33\text{ ms}$). This frame rate was chosen through experimental profiling. At 12~FPS, the time step is small enough to catch the sharp deceleration of a finger hitting the paper, while keeping the CPU load low so that MediaPipe tracking and PyTorch model inference run smoothly on ordinary CPUs without needing a GPU.

The MediaPipe Hand Landmarker is set for single-hand tracking (\texttt{num\_hands=1}). Focusing on one active hand improves CPU efficiency ($29.09\text{ ms}$ total processing latency) and prevents the hand from blocking the border AprilTags on compact A4 paper. The model provides 2D pixel coordinates for 21 hand landmarks $\{\mathbf{P}_i(t) = (x_i(t), y_i(t))\}_{i=0}^{20}$, as illustrated in Figure~\ref{fig:mediapipe_hand_skeleton}.

\begin{figure}[htbp]
\centering
\resizebox{0.70\textwidth}{!}{%
\begin{tikzpicture}[
    joint/.style = {circle, draw=blue!80!black, fill=blue!20!white, thick, inner sep=0pt, minimum size=7mm, font=\sffamily\bfseries\scriptsize},
    tip/.style   = {circle, draw=red!80!black, fill=red!20!white, thick, inner sep=0pt, minimum size=7mm, font=\sffamily\bfseries\scriptsize},
    wrist/.style = {circle, draw=purple!80!black, fill=purple!20!white, thick, inner sep=0pt, minimum size=7.5mm, font=\sffamily\bfseries\scriptsize},
    bone/.style  = {draw=gray!80!black, line width=1.5pt},
    normvec/.style = {draw=purple!80!black, dashed, line width=1.5pt, -{Stealth[length=3mm]}}
]

% Wrist
\node (0) [wrist] at (3.5, 0.0) {0};

% Thumb
\node (1) [joint] at (1.8, 1.2) {1};
\node (2) [joint] at (1.0, 2.2) {2};
\node (3) [joint] at (0.4, 3.2) {3};
\node (4) [tip]   at (0.0, 4.2) {4};

% Index
\node (5) [joint] at (2.4, 3.6) {5};
\node (6) [joint] at (2.2, 4.8) {6};
\node (7) [joint] at (2.0, 5.8) {7};
\node (8) [tip]   at (1.9, 6.8) {8};

% Middle
\node (9)  [joint] at (3.5, 3.8) {9};
\node (10) [joint] at (3.5, 5.2) {10};
\node (11) [joint] at (3.5, 6.4) {11};
\node (12) [tip]   at (3.5, 7.5) {12};

% Ring
\node (13) [joint] at (4.6, 3.6) {13};
\node (14) [joint] at (4.8, 4.9) {14};
\node (15) [joint] at (4.9, 6.0) {15};
\node (16) [tip]   at (5.0, 7.0) {16};

% Pinky
\node (17) [joint] at (5.6, 3.1) {17};
\node (18) [joint] at (6.0, 4.1) {18};
\node (19) [joint] at (6.3, 5.0) {19};
\node (20) [tip]   at (6.5, 5.9) {20};

% Bones
\draw [bone] (0) -- (1) -- (2) -- (3) -- (4);
\draw [bone] (0) -- (5) -- (6) -- (7) -- (8);
\draw [bone] (0) -- (9) -- (10) -- (11) -- (12);
\draw [bone] (0) -- (13) -- (14) -- (15) -- (16);
\draw [bone] (0) -- (17) -- (18) -- (19) -- (20);
\draw [bone] (5) -- (9) -- (13) -- (17);

% Normalization vector L_hand
\draw [normvec] (0) -- node[right=3pt, font=\sffamily\bfseries\footnotesize, text=purple!90!black] {$L_{\text{hand}} = \|\mathbf{P}_9 - \mathbf{P}_0\|_2$} (9);

% Labels
\node [above=2pt of 4, font=\sffamily\scriptsize\bfseries, text=red!80!black] {Thumb};
\node [above=2pt of 8, font=\sffamily\scriptsize\bfseries, text=red!80!black] {Index};
\node [above=2pt of 12, font=\sffamily\scriptsize\bfseries, text=red!80!black] {Middle};
\node [above=2pt of 16, font=\sffamily\scriptsize\bfseries, text=red!80!black] {Ring};
\node [above=2pt of 20, font=\sffamily\scriptsize\bfseries, text=red!80!black] {Pinky};
\node [below=2pt of 0, font=\sffamily\scriptsize\bfseries, text=purple!90!black] {Wrist Origin $\mathbf{P}_0$};

\end{tikzpicture}%
}
\caption{MediaPipe 21 hand landmarks topology and anatomical hand-length normalization reference ($L_{\text{hand}}$ between wrist $\mathbf{P}_0$ and middle MCP knuckle $\mathbf{P}_9$).}
\label{fig:mediapipe_hand_skeleton}
\end{figure}

\subsection{Unitless Hand-Length Scale Normalization}
\label{subsec:meth_scale_norm}
Because raw pixel coordinates change depending on how far the hand is from the camera lens and physical hand size, raw pixels cannot be fed directly to the classifier. To make the features independent of distance and hand size, coordinates are normalized using anatomical hand geometry.

The origin of the coordinate system is shifted to wrist landmark $\mathbf{P}_0(t) = (x_0(t), y_0(t))$. Hand length $L_{\text{hand}}(t)$ is defined as the Euclidean distance between the wrist and the middle metacarpophalangeal (MCP) knuckle (landmark 9):
\begin{equation}
L_{\text{hand}}(t) = \|\mathbf{P}_9(t) - \mathbf{P}_0(t)\|_2 = \sqrt{(x_9(t) - x_0(t))^2 + (y_9(t) - y_0(t))^2}
\label{eq:meth_hand_length_def}
\end{equation}

Unitless normalized coordinates $\mathbf{P}_i^{\text{norm}}(t)$ for each landmark $i \in \{0, \dots, 20\}$ are computed as:
\begin{equation}
\mathbf{P}_i^{\text{norm}}(t) = \left( \frac{x_i(t) - x_0(t)}{L_{\text{hand}}(t)}, \frac{y_i(t) - y_0(t)}{L_{\text{hand}}(t)} \right)
\label{eq:meth_norm_coords_def}
\end{equation}

Movement velocities $\mathbf{V}_i(t) = (v_{x,i}(t), v_{y,i}(t))$ between consecutive frames are derived using simple differences:
\begin{equation}
v_{x,i}(t) = \frac{x_i^{\text{norm}}(t) - x_i^{\text{norm}}(t-1)}{\Delta t}, \quad
v_{y,i}(t) = \frac{y_i^{\text{norm}}(t) - y_i^{\text{norm}}(t-1)}{\Delta t}
\label{eq:meth_velocity_def}
\end{equation}

\textbf{Direct Feature Propagation:} Normalized coordinates and velocities are passed directly into the neural network without applying temporal low-pass filters (such as Gaussian or moving average smoothing). Smoothing filters would flatten the sharp deceleration spike that signals the exact moment of physical paper impact.

\subsection{Temporal Sliding Window Construction}
\label{subsec:meth_sliding_window}
For each video frame $t$, the 21 normalized landmarks and their velocities form an 84-dimensional feature vector:
\begin{equation}
\mathbf{f}_t = \left[ x_0^{\text{norm}}, y_0^{\text{norm}}, v_{x,0}, v_{y,0}, \dots, x_{20}^{\text{norm}}, y_{20}^{\text{norm}}, v_{x,20}, v_{y,20} \right]^T \in \mathbb{R}^{84}
\label{eq:meth_feature_vector}
\end{equation}

Features are grouped into a temporal sliding window of \textbf{5 frames with a 2-frame overlap} (stride of 3 frames), producing an input sequence matrix:
\begin{equation}
\mathbf{X}_W = \left[ \mathbf{f}_{t-4}, \mathbf{f}_{t-3}, \mathbf{f}_{t-2}, \mathbf{f}_{t-1}, \mathbf{f}_t \right]^T \in \mathbb{R}^{5 \times 84}
\label{eq:meth_window_matrix}
\end{equation}
At 12~FPS, a 5-frame window covers roughly $416.7\text{ ms}$, which provides enough time to observe the full key strike motion: downward descent, surface deceleration, impact, and rebound.

\subsection{CPU Time Budget for 12~FPS Real-Time Processing}
\label{subsec:meth_timing_budget}
To maintain real-time performance on a commodity CPU, the processing time per frame $T_{\text{pipeline}}$ must stay comfortably below the 12~FPS budget:
\begin{equation}
T_{\text{budget}} = \frac{1}{12\text{ s}} \approx 83.33\text{ ms}
\label{eq:meth_timing_budget}
\end{equation}

The average time taken by each component on a standard quad-core CPU is:
\begin{itemize}
    \item Video frame acquisition: $T_{\text{capture}} \approx 3.25\text{ ms}$
    \item AprilTag corner tracking and homography: $T_{\text{apriltag}} \approx 6.84\text{ ms}$
    \item MediaPipe 21 hand joint tracking: $T_{\text{mediapipe}} \approx 7.92\text{ ms}$
    \item Scale normalization and velocity derivation: $T_{\text{norm}} \approx 0.45\text{ ms}$
    \item PyTorch LSTM model forward pass: $T_{\text{lstm}} \approx 2.08\text{ ms}$
    \item Coordinate mapping and key dispatch: $T_{\text{dispatch}} \approx 8.55\text{ ms}$
\end{itemize}

Summing these times yields a total end-to-end latency of $T_{\text{pipeline}} \approx 29.09\text{ ms}$ ($\approx 34.4\text{ FPS}$). This leaves a safety margin of $54.24\text{ ms}$ (over $65\%$ margin) within the $83.33\text{ ms}$ window, ensuring that the software runs without lagging or dropping frames on standard consumer computers.

\section{Deep Learning Sequence Model Architecture}
\label{sec:meth_dl_models}

\subsection{Multi-Finger Touch Contact Classification}
\label{subsec:meth_touch_formulation}
Touch detection is treated as a multi-label classification problem. For each 5-frame window $\mathbf{X}_W$, the neural network predicts independent contact probabilities for all five fingers of the active hand:
\begin{equation}
\hat{\mathbf{y}} = \left[ \hat{y}_{\text{thumb}}, \hat{y}_{\text{index}}, \hat{y}_{\text{middle}}, \hat{y}_{\text{ring}}, \hat{y}_{\text{pinky}} \right] \in [0, 1]^5
\label{eq:meth_multilabel_prob}
\end{equation}
Evaluating all five digits in parallel removes the single-finger limitation of older shadow-based systems, allowing natural typing with any finger.

\subsection{PyTorch LSTM Sequence Classifier}
\label{subsec:meth_lstm_arch}
To achieve real-time inference on a CPU without a GPU, a lightweight uni-directional Long Short-Term Memory (LSTM) network was developed in PyTorch. The model structure consists of:
\begin{itemize}
    \item \textbf{Input Layer:} Takes the sequence matrix $\mathbf{X}_W \in \mathbb{R}^{B \times 5 \times 84}$, where $B$ is batch size.
    \item \textbf{LSTM Hidden Layers:} Two stacked uni-directional LSTM layers with 64 hidden units and dropout ($p = 0.20$):
    \begin{equation}
    \mathbf{h}_t, \mathbf{c}_t = \text{LSTM}(\mathbf{f}_t, \mathbf{h}_{t-1}, \mathbf{c}_{t-1})
    \label{eq:meth_lstm_step}
    \end{equation}
    \item \textbf{Dense Linear Layer:} A fully connected layer projecting the final hidden state $\mathbf{h}_5 \in \mathbb{R}^{64}$ to 5 output logits.
    \item \textbf{Sigmoid Output:} Computes independent touch probabilities $\hat{y}_k = \sigma(z_k) = \frac{1}{1 + e^{-z_k}}$.
\end{itemize}

The network is trained using the multi-label Binary Cross-Entropy loss ($\mathcal{L}_{\text{BCE}}$):
\begin{equation}
\mathcal{L}_{\text{BCE}} = -\frac{1}{5} \sum_{k=1}^{5} \left[ y_k \log(\hat{y}_k) + (1 - y_k) \log(1 - \hat{y}_k) \right]
\label{eq:meth_bce_loss}
\end{equation}
A touch event is confirmed when the predicted probability for digit $k$ crosses the threshold $\tau_{\text{touch}} = 0.90$.

\section{Application 2: Runtime Virtual Keyboard Engine}
\label{sec:meth_runtime_engine}
Application 2 serves as the runtime typing software, divided into an interactive setup GUI (Component A) and a background vision watch engine (Component B).

\subsection{Planar Homography Derivation ($H$)}
\label{subsec:meth_homography_derivation}
To map fingertip pixel coordinates $(u, v)$ from the camera image to the paper layout coordinates $(X, Y)$, the system computes a $3 \times 3$ Planar Homography matrix $H$:
\begin{equation}
\begin{bmatrix} X \\ Y \\ 1 \end{bmatrix} \sim H \begin{bmatrix} u \\ v \\ 1 \end{bmatrix} = \begin{bmatrix} h_{11} & h_{12} & h_{13} \\ h_{21} & h_{22} & h_{23} \\ h_{31} & h_{32} & h_{33} \end{bmatrix} \begin{bmatrix} u \\ v \\ 1 \end{bmatrix}
\label{eq:meth_homography_matrix}
\end{equation}

The AprilTag detector finds the corner points of all visible border markers with sub-pixel precision. Given $N \ge 4$ point pairs between image pixels $(u_j, v_j)$ and known XML paper coordinates $(X_j, Y_j)$, a system of linear equations $\mathbf{A} \mathbf{h} = \mathbf{0}$ is formed:
\begin{equation}
\begin{bmatrix}
-u_j & -v_j & -1 & 0 & 0 & 0 & u_j X_j & v_j X_j & X_j \\
0 & 0 & 0 & -u_j & -v_j & -1 & u_j Y_j & v_j Y_j & Y_j
\end{bmatrix} \mathbf{h} = \mathbf{0}
\label{eq:meth_dlt_matrix}
\end{equation}
The 9-element vector $\mathbf{h}$ is solved using Singular Value Decomposition (SVD), taking the column vector corresponding to the smallest singular value of matrix $\mathbf{A}$.

Because $H$ is updated continuously whenever the paper or camera moves, the virtual keyboard handles perspective tilt up to $75^\circ$, camera vibration, and paper movement without requiring a rigid mount or manual recalibration.

\subsection{Fingertip Coordinate Transformation and Key Lookup}
\label{subsec:meth_coord_transform}
When the LSTM model detects a touch from finger $k$, the 2D image coordinates of that fingertip landmark $(u_k, v_k)$ (landmark 4, 8, 12, 16, or 20) are projected into XML layout space:
\begin{equation}
x' = h_{11} u_k + h_{12} v_k + h_{13}, \quad y' = h_{21} u_k + h_{22} v_k + h_{23}, \quad w' = h_{31} u_k + h_{32} v_k + h_{33}
\label{eq:meth_proj_steps}
\end{equation}
\begin{equation}
X_{\text{XML}} = \frac{x'}{w'}, \quad Y_{\text{XML}} = \frac{y'}{w'}
\label{eq:meth_xml_coords}
\end{equation}

The mapped coordinate $(X_{\text{XML}}, Y_{\text{XML}})$ is compared against the bounding boxes of all keys in the loaded XML file. If $x_{\min}^{(m)} \le X_{\text{XML}} \le x_{\max}^{(m)}$ and $y_{\min}^{(m)} \le Y_{\text{XML}} \le y_{\max}^{(m)}$, key $m$ is registered as pressed.

\subsection{Action Multiplexing and Command Dispatch}
\label{subsec:meth_action_multiplexing}
The system separates physical layout geometry from digital actions. In Application 2 Component A, users can link any key ID to either:
\begin{itemize}
    \item \textbf{Keystroke Action:} A standard OS keyboard input (such as letters, numbers, spaces, or modifier keys).
    \item \textbf{System Shell Command Action:} A computer command or script (such as opening a terminal program, controlling audio volume, or launching custom Python scripts).
\end{itemize}

Action mappings are saved in JSON configuration files. Users can easily change configuration profiles (for example, switching from standard typing mode to developer shortcut mode or gaming mode) on the exact same printed sheet of paper without needing to print another page.

\section{Chapter Summary}
\label{sec:meth_summary}
This chapter explained the research methodology, system design, and algorithmic steps of the paper-based virtual keyboard. Following a positivist approach and the DSRM framework, the system combines a PySide6 Layout Designer with a multi-threaded Runtime Engine. Key technical components include: 12~FPS video resampling, unitless hand scale normalization ($L_{\text{hand}}$) with direct feature propagation, multi-finger PyTorch LSTM touch detection running on standard CPUs, and dynamic AprilTag planar homography tracking ($H$). Chapter~\ref{ch:results} presents the experimental benchmarking results and performance tests across the system.
